\begin{corollary}[6.12.6]
\label{IV.2.s6.12.6-fr}
Sea $A$ un anillo que posee una de las propiedades siguientes:
\begin{enumerate}
\item[\textup{(i)}] $A$ es un anillo de Dedekind y su cuerpo de fracciones $K$ es
de característica \textup{0}.
\item[\textup{(ii)}] $A$ es un anillo semilocal noetheriano de dimensión $\leq1$.
\end{enumerate}
Entonces, para todo preesquema $X$ localmente de tipo finito sobre $A$,
$\operatorname{Reg}(X)$ es abierto en $X$.
\end{corollary}

\begin{proof}
Verifiquemos en los dos casos la condición \emph{c)} de (6.12.4). En ambos casos,
un ideal primo $\mathfrak p$ de $A$ es maximal o minimal; si
$\mathfrak p$ es maximal, una $(A/\mathfrak p)$-álgebra finita e íntegra
es un cuerpo y la condición \emph{c)} de (6.12.4) se verifica trivialmente.
Supongamos, pues, que $\mathfrak p$ no es maximal y distingamos los dos casos
del enunciado.

(i) Si $K$ es de característica 0, no existe otra extensión radicial de
$K$ que el propio $K$; como un anillo de Dedekind es regular (\textbf 0, 17.1.4),
la condición \emph{c)} de (6.12.4) se verifica trivialmente.

(ii) Se puede suponer entonces que $A$ es íntegro (6.12.2); sea $K$ su cuerpo
de fracciones. Si $K'$ es una extensión radicial finita de $K$ y $A'$ una
sub-$A$-álgebra de $K'$ engendrada por un sistema finito de generadores de
$K'$ sobre $K$, enteros sobre $A$, entonces $A'$ es un anillo semilocal
íntegro de dimensión 1 (\textbf 0, 16.1.5); por tanto, en
$X'=\operatorname{Spec}(A')$, el conjunto formado solo por el punto genérico
es abierto y está evidentemente contenido en $\operatorname{Reg}(X')$, lo
que demuestra la condición \emph{c)} de (6.12.4) en este caso.
\end{proof}

Este corolario se aplica, en particular, cuando $A=\mathbf Z$.
